\clearpage
\section{Conclusion}
\label{sec:Conclusion}
\label{sec:conclusion}

Agentic systems take over multi-step work in organizations, but cost and reliability still limit their use. For every task of such a system, a software architect must decide whether predefined code or the model controls the progression of the task. Practice has made this choice by experience and rules of thumb, and no reviewed approach combined task characteristics, comparative evidence, constraints and priorities to support it. This thesis therefore asked how a task-level architecture decision framework can support this choice. Following design science research, it built TADF on paired implementations of both orchestration paradigms and revised it through practitioner interviews, a benchmark evaluation and a demonstration.

The paired comparison shows that neither orchestration paradigm is better in general. The agent pays off where the next step or the sources cannot be fixed before the run. Where later actions depended on earlier outcomes, it solved 13 of 15 executions and the LLM workflow none. Where written rules decided, the LLM workflow was more reliable, and in six of eight task archetypes the agent needed 1.8 to 7.7 times the tokens per successful execution. The choice is therefore a sequence of levels rather than one decision. For current models and production systems, the evidence suggests the LLM workflow as the default and the agent as a bounded component for the parts that cannot be fixed in advance.

The answer to the research question is a framework that follows these levels. TADF qualifies a task within its system context, checks whether its control requirements call for an agent and then weighs the priorities of the architect, within hard limits that no weight can offset. Each recommendation names its evidence and adds the design knowledge for building it, including a mandatory minimum for every LLM workflow. In its final form for this thesis, TADF is a skill for an LLM. The model interprets the description, code applies the fixed rules and the architect confirms. The evaluation supports this design in part. Several practitioners favored such support inside their LLM assistants and the current discurse also point ito that direction. Because code applies the rules, the same answers lead to the same recommendation. On WorkBench, however, routing by TADF only equaled Always LLM workflow, and no independent architect has yet decided with TADF. The hypothesis of consistent and practically useful recommendations therefore holds only in part.

The thesis thus contributes to knowledge and practice. Its decision method and four design principles turn the comparative evidence into design knowledge: decide from coarse to fine, state the implementation behind every recommendation, provide the consultation where architects already work, and let code apply the rules. In the terms of \textcite{gregorPositioningPresentingDesign2013}, this is an improvement: an explicit, evidence-linked procedure for a known decision that practice has so far addressed with heuristics. For practice, software architects gain a consultation in the assistant they already use. Clients and managers gain a simple question to ask before they pay for an agent: does the task need one?

The comparative evidence rests on implementations built for this thesis, three OpenAI models and one external benchmark, and the builder of TADF also took the architect role. Because models change quickly, TADF is best kept as a living framework whose rules are revised one by one as new evidence arrives. The next step is a study in which independent architects decide comparable design cases with TADF, with their usual practice and with the same assistant without TADF. These limits do not remove the value of the work. In a young field, the thesis provides paired evidence at the task level, a working instrument and design principles that future research can test, extend and revise.

The question in the title, LLM workflow or agent, therefore has no general answer. In line with the hypothesis it has an answer per task. Parts of the current discourse treat the agent as the goal of automation. This thesis places it as one option among others. The agent is worth its cost where the model must choose the next step, and predefined control is usually the better choice where the path can be written down in advance. For research on agentic information systems and decision support, the results move attention from what agents can do to how architects decide where they should act. As models grow stronger and architectures evolve, this decision will not disappear. It will move, and decision support has to move with it.




